\chapter{Nation, Civilization, and Collective Myth}
\label{ch:nation-myth}

\epigraph{Nations are not something eternal. They had their beginnings and they will end.}{Ernest Renan, \textit{What Is a Nation?} (1882)}

\noindent
National narratives are among the most expensive doxonomic products in history.
Every modern state invests in the production and maintenance of national identity: through education, monuments, ceremonies, media, and military ritual.

\section{National Narrative as Doxonomic Product}
\label{sec:ch27-national}

\begin{definition}[National Myth]
\label{def:national-myth}
\index{national myth}
A \emph{national myth} is a structured narrative $n_{\text{nat}} \in \narspace$ specifying:
\begin{enumerate}[nosep]
  \item a founding story (origin, sacrifice, triumph),
  \item a collective character (virtues attributed to the nation),
  \item an enemy or other (defining the boundary of ``us''),
  \item a destiny (what the nation is called to achieve).
\end{enumerate}
\end{definition}

\begin{model}[Myth Maintenance]
\label{mod:myth-maintenance}
\index{myth maintenance}
The strength of national identification $\belief{N}$ requires continuous investment:
\[
  \frac{d\belief{N}}{dt} = \alpha \phi_N - \delta \belief{N} + \beta \belief{N}(1 - \belief{N}) - \gamma \cdot D(t)
\]
where $\phi_N$ is the institutional investment, $\delta$ is natural decay, $\beta$ is social reinforcement, and $D(t)$ measures the dissonance from events contradicting the myth.
\end{model}

\section{Historical Memory Management}
\label{sec:ch27-memory}

States actively manage historical memory: what is remembered, how it is framed, and what is omitted.
Curriculum design, museum curation, memorial construction, and holiday designation are all doxonomic operations with measurable budgets \autocite{anderson2006}.

\section{Notes and References}
\label{sec:ch27-notes}

On nations as imagined communities, see \textcite{anderson2006}.
State simplification and legibility are analyzed in \textcite{scott1998}.
On propaganda and national myth, see \textcite{lippmann1922} and \textcite{bernays1928}.
For the political economy of nationalism, see \textcite{acemoglu2012}.

\section{Exercises}

\begin{exercise}
Estimate the annual doxonomic investment in national identity for a country of your choice.
Include education, military ceremony, state media, monuments, and holidays.
\end{exercise}

\begin{exercise}
Using the myth maintenance model, predict the effect of a major scandal on national identification.
At what $D(t)$ does the myth become unsustainable without additional investment?
\end{exercise}

\begin{exercise}
Compare the national myths of two countries.
Identify the structural parallels in their belief production chains.
\end{exercise}
